% Generated draft: requires mathematical and source review.
\documentclass[11pt]{article}
\usepackage[margin=25mm]{geometry}
\usepackage{fontspec}
\setmainfont{Latin Modern Roman}
\usepackage{amsmath,amssymb,mathtools,unicode-math}
\setmathfont{Latin Modern Math}
\usepackage{enumitem,needspace,xurl,hyperref}
\setlength{\parindent}{0pt}
\setlength{\parskip}{5pt}
\setlength{\emergencystretch}{3em}
\newcommand{\N}{\mathbb{N}}
\newcommand{\Z}{\mathbb{Z}}
\newcommand{\Q}{\mathbb{Q}}
\newcommand{\R}{\mathbb{R}}
\newcommand{\C}{\mathbb{C}}
\newcommand{\F}{\mathbb{F}}
\newcommand{\A}{\mathbb{A}}
\newcommand{\PP}{\mathbb{P}}
\newcommand{\E}{\mathbb{E}}
\newcommand{\eps}{\varepsilon}
\newcommand{\dd}{\,\mathrm{d}}
\begin{document}
\section*{UnsolvedMath ID: 30003874}
\textbf{Source problem code:} catalogue entry
\subsection{Definitions and mathematical statement}
For a specified ground field \(k\), let \(f:Y\to X\) be a dominant generically finite map of irreducible \(k\)-varieties. Its resolvent degree \(\operatorname{RD}_k(f)\) is the least nonnegative integer \(d\) for which there is a tower of dominant generically finite rational maps
\[
X_r\dashrightarrow X_{r-1}\dashrightarrow\cdots\dashrightarrow X_0=X
\]
and a dominant rational map \(X_r\dashrightarrow Y\) over \(X\), such that each step of the tower is, birationally, obtained by base change from a dominant generically finite map of varieties of dimension at most \(d\). This definition concerns dense open subsets; it does not prescribe an extension across degenerate fibers. The supplied record does not specify the ground field or precise models of its parameter spaces, so these data remain unresolved rather than being chosen here.

The notation \(H_{3,3}\) denotes the source's space parametrizing smooth cubic surfaces, and \(H_{3,3}^{(1)}\) denotes its incidence cover parametrizing pairs \((S,L)\), where \(S\) is such a surface and \(L\) is a line contained in \(S\). The map forgets \(L\) and is described in the source as a \(27\)-sheeted cover. Likewise, \(H_{4,2}\) denotes the source's space parametrizing smooth plane quartics, and \(H_{4,2}^{(1)}\) denotes its incidence cover parametrizing pairs \((C,B)\), where \(B\) is a bitangent line to \(C\). This map forgets \(B\) and is described as a \(28\)-sheeted cover. On the ordinary-bitangent locus, this means that the intersection divisor is \(B\cdot C=2p+2q\) for distinct geometric points \(p,q\). The record does not specify its convention for hyperflexes or the precise locus on which its cover description applies.

With these unresolved choices understood, the supplied conjectural targets are
\[
\operatorname{RD}\!\left(H_{3,3}^{(1)}\longrightarrow H_{3,3}\right)=3
\]
and
\[
\operatorname{RD}\!\left(H_{4,2}^{(1)}\longrightarrow H_{4,2}\right)
=\dim H_{4,2}=6.
\]
Here \(\operatorname{RD}\) suppresses the unspecified ground field. The asserted dimension \(6\) is retained from the record; the record does not define whether these spaces are moduli spaces, stacks, or other parameter spaces. Consequently, this is a qualified transcription of the conjectures, pending specification of those data. The background literature assessment refers instead to genus-two curves, which is inconsistent with the stated smooth-plane-quartic target: smooth plane quartics have genus three. No replacement of either target is made.
\subsection{Short English statement}
The supplied conjectures assign resolvent degree \(3\) to choosing a line on a smooth cubic surface and resolvent degree \(6\) to choosing a bitangent of a smooth plane quartic. The ground field, parameter-space models, and bitangent boundary conventions require clarification. The supplied background's reference to genus-two curves conflicts with the quartic statement. These are statement-only conjectural drafts requiring mathematical review.
\subsection{Sources}
\begin{itemize}
\item Local UnsolvedMath registry, supplied complete source record: permanent ID \(30003874\), title ``Resolvent Degrees of Lines and Bitangents,'' source problem code \texttt{catalogue entry}, collection ``Oberwolfach Reports Open Problems.'' No individual proposer is supplied; the proposed year is recorded as \(2018\).
\item Supplied report citation: \emph{Algebraische Zahlentheorie (2019)}, Oberwolfach Reports, pp.\ 1785--1844, DOI \texttt{10.4171/owr/2018/30}. \url{https://doi.org/10.4171/owr/2018/30}. The citation's year \(2019\) and the record's \(2018\) year information are retained without reconciliation.
\item Supplied literature source: ``Farb--Wolfson, resolvent degree.'' \url{https://arxiv.org/abs/1803.04063}.
\item Supplied literature source: ``Recent resolvent-degree study.'' \url{https://escholarship.org/content/qt0598s939/qt0598s939.pdf}.
\item Literature assessment dated \texttt{2026-08-21}: the supplied curation labels the problem open and says that the conjectural resolvent degrees for cubic surfaces and genus-two curves are unproved, with partial geometric interpretations and bounds available. This is supplied dated curation, not an independently verified assessment. Its genus-two wording conflicts with the plane-quartic formulation. The record itself asks that time-sensitive status and cited claims be verified independently.
\item The supplied statement status is \texttt{corrected\_verified}, with a reported manual semantic, typographic, and LaTeX review. This records the supplied editorial assessment and does not establish mathematical correctness or independent source verification.
\end{itemize}
\end{document}
